\chapter{Unit 1: Introduction to Engineering Drawing \& AutoCAD Setup}

\section{A. What You'll Learn}
Engineering drawing is the universal graphic language used by engineers and architects to communicate physical object designs with mathematical precision. In this unit, you will master:
\begin{itemize}
    \item Essential manual drawing instruments, pencil lead grading, and proper handling.
    \item Standardized Line Types (The Alphabet of Lines) defined by ISO and ANSI standards.
    \item Dimensioning components, mechanics, Aligned vs. Unidirectional systems, and single-stroke vertical Gothic lettering.
    \item Scales, Representative Fraction ($R.F.$), Length of Scale ($L.O.S.$), Plain Scales, and Diagonal Scales based on similar triangles.
    \item Fundamental AutoCAD 2D interface navigation, drawing limits, units, object snap (\texttt{OSNAP}), orthogonal mode (\texttt{ORTHO}), and function keys.
\end{itemize}

---

\section{B. Conceptual Framework of Drawing Instruments}

\subsection{Key Instruments and Their Applications}

\begin{definition}[Drawing Board]
Provides a smooth, flat, rectangular surface for securing drawing sheets. Made of seasoned soft wood strips (such as pine) to prevent warping and allow easy pin/tape insertion.
\end{definition}

\begin{definition}[Mini-Drafter]
An indispensable drafting tool that combines the functions of a T-square, set squares, protractor, and scales into a single mechanism. It uses a 3D parallelogram linkage clamped to the drawing board to keep ruling edges parallel across the entire sheet.
\end{definition}

\begin{itemize}
    \item \textbf{T-Square:} Traditional tool consisting of a stock and blade used primarily for drawing horizontal parallel lines and supporting set squares.
    \item \textbf{Set Squares:} Used in combination with a drafter or T-square:
    \begin{itemize}
        \item \textbf{$45^\circ\text{ Set Square}$:} Angles of $45^\circ, 90^\circ, 135^\circ$.
        \item \textbf{$30^\circ\text{--}60^\circ\text{ Set Square}$:} Angles of $30^\circ, 60^\circ, 90^\circ, 120^\circ, 150^\circ$.
    \end{itemize}
    \item \textbf{Instrument Box:}
    \begin{itemize}
        \item \textbf{Large Compass:} Features a knee joint for drawing large circles ($> 50\text{ mm}$ radius).
        \item \textbf{Bow Compass:} Used for drawing precise small circles and arcs.
        \item \textbf{Dividers:} Used to transfer dimensions from scales to drawing paper or divide lines into equal segments.
    \end{itemize}
\end{itemize}

\subsection{Pencil Lead Grading and Selection}
Pencil leads are classified by hardness (clay content vs. graphite ratio).

\begin{importantrule}[Pencil Lead Application Matrix]
\begin{itemize}
    \item \textbf{Hard Grades ($4\text{H to } 6\text{H}$):} High clay content. Used for very light construction lines, center lines, locus lines, and guide lines.
    \item \textbf{Medium Grades ($\text{H}, 2\text{H}$):} Standard drafting leads. Used for visible object lines, dimension lines, extension lines, arrowheads, and lettering.
    \item \textbf{Soft Grades ($\text{HB}, \text{B}$):} High graphite content. Used for artistic sketching or freehand boundary lines. (Rarely used in precision engineering drawing).
\end{itemize}
\end{importantrule}

\subsection{French Curves}
French curves are plastic templates with varying smooth curves used to draw non-circular curves (such as ellipses, parabolas, hyperbolas, and involutes) that cannot be drawn with a standard compass.

---

\section{C. Line Types (The Alphabet of Lines)}

Standardization of lines is governed by ISO and ANSI standards. Line thickness and style convey specific geometric meanings.

\begin{table}[htbp]
\centering
\small
\begin{tabularx}{\textwidth}{l c X c}
\toprule
\textbf{Line Type} & \textbf{Representation} & \textbf{Description \& Application} & \textbf{Pencil Grade} \\
\midrule
\textbf{Visible / Object Line} & \rule{2.5cm}{1.2pt} & Continuous thick line. Represents visible outlines and sharp edges. & $\text{H}$ or $\text{HB}$ \\
\addlinespace
\textbf{Hidden Line} & \rule[0.5ex]{0.4cm}{0.8pt}\hspace{2pt}\rule[0.5ex]{0.4cm}{0.8pt}\hspace{2pt}\rule[0.5ex]{0.4cm}{0.8pt} & Dashed medium line. Represents interior/hidden edges. & $2\text{H}$ \\
\addlinespace
\textbf{Center Line} & \rule[0.5ex]{0.6cm}{0.5pt}\hspace{2pt}\rule[0.5ex]{0.1cm}{0.5pt}\hspace{2pt}\rule[0.5ex]{0.6cm}{0.5pt} & Long dash -- short dash chain line. Indicates axes of symmetry. & $4\text{H}$ \\
\addlinespace
\textbf{Construction Line} & \rule{2.5cm}{0.3pt} & Continuous very thin line. Used for initial layout projections. & $4\text{H}$ or $6\text{H}$ \\
\addlinespace
\textbf{Dimension / Extension} & \rule{2.5cm}{0.5pt} & Continuous thin line. Defines start and limit of measurements. & $2\text{H}$ \\
\addlinespace
\textbf{Cutting Plane Line} & \rule[0.5ex]{0.6cm}{1.5pt}\hspace{2pt}\rule[0.5ex]{0.1cm}{1.5pt}\hspace{2pt}\rule[0.5ex]{0.6cm}{1.5pt} & Thick chain line with thickened ends. Shows hypothetical section cut. & $\text{H}$ \\
\bottomrule
\end{tabularx}
\caption{The Alphabet of Lines: Standard Engineering Line Types.}
\label{tab:line-types}
\end{table}

---

\section{D. Dimensioning Principles \& Systems}

Dimensioning is the process of specifying exact physical sizes, locations, and geometric characteristics of object features.

\subsection{Components of a Dimension}
\begin{itemize}
    \item \textbf{Extension Line:} Thin continuous line extending from the object edge. A gap of $1\text{ mm}$ is left between object and extension line.
    \item \textbf{Dimension Line:} Thin continuous line terminated by arrowheads defining measurement boundaries.
    \item \textbf{Arrowheads:} Standard ratio of **Length to Width is $3:1$**. Arrowheads must be sharp, filled, and uniform.
    \item \textbf{Leader Line:} Thin line leading from a note to a feature (e.g., hole diameter), terminating in an arrowhead or dot.
\end{itemize}

\subsection{Systems of Dimensioning}

\begin{definition}[Aligned System]
Dimensions are placed perpendicular to the dimension line. Values are placed above the dimension line and read from the bottom or right-hand side of the sheet.
\end{definition}

\begin{definition}[Unidirectional System]
Dimensions are placed vertically regardless of line orientation. Values are read from the bottom of the sheet only. The dimension line is broken in the middle to insert text.
\end{definition}

\begin{table}[htbp]
\centering
\small
\begin{tabularx}{\textwidth}{l X X}
\toprule
\textbf{Feature} & \textbf{Aligned Dimensioning System} & \textbf{Unidirectional Dimensioning System} \\
\midrule
\textbf{Text Orientation} & Perpendicular to dimension line & Always horizontal (upright) \\
\textbf{Reading Direction} & Read from bottom or right side & Read from bottom only \\
\textbf{Dimension Line} & Continuous unbroken line & Broken in middle for text entry \\
\textbf{Application} & Standard engineering drawings (ISO) & Aircraft, automotive, heavy machinery (ANSI) \\
\bottomrule
\end{tabularx}
\caption{Comparison of Aligned and Unidirectional Dimensioning Systems.}
\label{tab:dimension-systems}
\end{table}

---

\section{E. Single-Stroke Vertical Gothic Lettering}

Lettering on technical drawings must be uniform, highly legible, and suitable for microfilming or digital scanning. "Single stroke" means letter line thickness is produced in one stroke of the pencil.

\begin{keyequation}[Single-Stroke Lettering Proportions]
\begin{itemize}
    \item \textbf{Standard Height-to-Width Ratio:} $7:4$ for most letters ($A, B, C, D, E, F, G, H, J, K, L, N, O, P, Q, R, S, T, U, V, X, Y, Z$).
    \item \textbf{Wide Letters ($M, W$):} Ratio of $7:5$ or $7:6$.
    \item \textbf{Narrow Letter ($I$):} Single vertical stroke ($7:1$).
    \item \textbf{Letter Spacing:} $\approx \frac{1}{5} h$ (where $h$ is character height).
    \item \textbf{Word Spacing:} $\approx 1 h$.
    \item \textbf{Line Spacing:} $\approx \frac{1}{2} h \text{ to } 1 h$.
\end{itemize}
\end{keyequation}

---

\section{F. Scales and Representative Fraction ($R.F.$)}

Scales allow large objects (e.g., buildings, bridges) or small objects (e.g., watch gears, micro-chips) to be represented accurately on standard paper sizes.

\begin{keyequation}[Representative Fraction ($R.F.$)]
\begin{equation}
    R.F. = \frac{\text{Length of object on drawing sheet}}{\text{Actual physical length of object}}
\end{equation}
\end{keyequation}

\begin{keyequation}[Length of Scale ($L.O.S.$)]
\begin{equation}
    L.O.S. = R.F. \times \text{Maximum length to be measured}
\end{equation}
\end{keyequation}

\subsection{Classification of Scales}
\begin{itemize}
    \item \textbf{Full Scale ($1:1$):} Drawing dimensions equal actual physical dimensions.
    \item \textbf{Reducing Scale ($1:X$):} Drawing dimensions smaller than actual object (e.g., $1:10, 1:50, 1:100$).
    \item \textbf{Enlarging Scale ($X:1$):} Drawing dimensions larger than actual object (e.g., $2:1, 5:1, 10:1$).
\end{itemize}

\subsection{Plain Scale vs. Diagonal Scale}

\begin{definition}[Plain Scale]
Consists of a line divided into a number of equal main units, with the first main unit subdivided into smaller parts. Represents two units (e.g., Kilometers and Hectometers, or Meters and Decimeters).
\end{definition}

\begin{definition}[Diagonal Scale]
Used when minute measurements are required or to measure three consecutive units (e.g., Meters, Decimeters, and Centimeters) or up to two decimal places. Based on the principle of \textbf{Similar Triangles}.
\end{definition}

\begin{importantrule}[Principle of Similar Triangles in Diagonal Scales]
To divide a short vertical line of length $h$ into 10 equal parts, a diagonal line is drawn across 10 horizontal parallel lines. By similar triangles, the distance from the vertical axis to the diagonal line at the $n^{\text{th}}$ horizontal line equals $\frac{n}{10}$ of the base unit.
\end{importantrule}

\begin{example}[Diagonal Scale Construction]
\textbf{Problem:} Construct a diagonal scale of $R.F. = 1:40$ to read meters, decimeters, and centimeters up to $6\text{ m}$. Mark a distance of $4.56\text{ m}$ on the scale.\\[4pt]
\textbf{Solution Step 1: Calculate Length of Scale ($L.O.S.$)}
\begin{equation*}
    L.O.S. = R.F. \times \text{Max Length} = \frac{1}{40} \times 6\text{ m} = \frac{1}{40} \times 600\text{ cm} = \mathbf{15\text{ cm}}
\end{equation*}

\noindent \textbf{Solution Step 2: Main Division}
Draw a line $15\text{ cm}$ long and divide it into 6 equal main parts. Each main part represents $1\text{ m}$.

\noindent \textbf{Solution Step 3: Subdivisions \& Diagonals}
Divide the first left division ($1\text{ m}$) horizontally into 10 equal parts (each $= 1\text{ dm} = 0.1\text{ m}$). Erect a vertical rectangle height $5\text{ cm}$ and divide vertically into 10 equal parts (each $= 1\text{ cm} = 0.01\text{ m}$). Draw diagonal lines connecting subdivision $n$ on bottom to $n+1$ on top.

\noindent \textbf{Solution Step 4: Locate $4.56\text{ m}$}
Take 4 main units ($4\text{ m}$) to the right of zero, move 5 subdivisions left ($0.5\text{ m}$), and move up to the $6^{\text{th}}$ horizontal line ($0.06\text{ m}$). Total distance = $4 + 0.5 + 0.06 = \mathbf{4.56\text{ m}}$.
\end{example}

---

\section{G. Introduction to AutoCAD Interface \& Setup}

AutoCAD (Computer-Aided Design) replaces manual drafting tools with digital precision.

\subsection{Interface Layout Components}
\begin{enumerate}
    \item \textbf{Application Menu (Big Red A):} File management (New, Open, Save, Export).
    \item \textbf{Ribbon:} Organized tabs (Home, Insert, Annotate) containing panels of drafting tools.
    \item \textbf{Drawing Area (Viewport):} Infinite 2D/3D graphics area.
    \item \textbf{UCS Icon:} Visual indicator of $X, Y, Z$ coordinate axes.
    \item \textbf{Command Line:} Primary method for entering AutoCAD commands and viewing system prompts.
    \item \textbf{Status Bar:} Toggles for drawing aids (Grid, Snap, Ortho, OSNAP) at bottom right.
\end{enumerate}

\subsection{Essential AutoCAD Setup Commands}

\begin{autocadcmd}[UNITS]
Defines measurement type and precision.
\begin{itemize}
    \item \textbf{Command:} \texttt{UNITS} or \texttt{UN}
    \item \textbf{Type:} Decimal (Mechanical), Architectural (Feet/Inches).
    \item \textbf{Precision:} Number of decimal places (e.g., \texttt{0.00}).
    \item \textbf{Insertion Scale:} Units for blocks/content (Millimeters).
\end{itemize}
\end{autocadcmd}

\begin{autocadcmd}[LIMITS]
Sets the invisible boundary of the drawing area (equivalent to paper size).
\begin{itemize}
    \item \textbf{Command:} \texttt{LIMITS}
    \item \textbf{Lower-Left Corner:} Type \texttt{0,0} $\to$ Press \texttt{Enter}.
    \item \textbf{Upper-Right Corner:} Type \texttt{210,297} (for A4 sheet) $\to$ Press \texttt{Enter}.
    \item \textbf{Required Refresh:} Type \texttt{Z} (Zoom) $\to$ \texttt{Enter} $\to$ Type \texttt{A} (All) $\to$ \texttt{Enter}.
\end{itemize}
\end{autocadcmd}

\subsection{Function Key Cheat Sheet}
\begin{table}[htbp]
\centering
\small
\begin{tabularx}{\textwidth}{c l X}
\toprule
\textbf{Key} & \textbf{Function} & \textbf{Operational Description} \\
\midrule
\texttt{F1} & Help & Opens Autodesk help documentation. \\
\texttt{F2} & Text Screen & Displays command history window. \\
\texttt{F3} & OSNAP & Toggles Object Snap on/off (Endpoint, Midpoint, Center). \\
\texttt{F7} & Grid & Toggles background grid display. \\
\texttt{F8} & ORTHO & Restricts cursor motion to horizontal ($0^\circ, 180^\circ$) and vertical ($90^\circ, 270^\circ$). \\
\texttt{F9} & Snap Mode & Restricts cursor movement to grid intervals. \\
\texttt{F10} & Polar Tracking & Restricts cursor to specified angle increments. \\
\texttt{F12} & Dynamic Input & Displays command prompts and coordinates near cursor. \\
\bottomrule
\end{tabularx}
\caption{AutoCAD Function Key Reference.}
\label{tab:function-keys}
\end{table}

---

\section{H. Chapter 1 Practice Problems \& MCQs}

\subsection{Practice Questions}

\begin{vivaqa}[Question 1 (Basic)]
State the standard height-to-width ratio for single-stroke vertical Gothic lettering for normal letters, wide letters ($M, W$), and narrow letter ($I$).\\[4pt]
\textbf{Answer:} Normal letters $= 7:4$; Wide letters ($M, W$) $= 7:5$ or $7:6$; Narrow letter ($I$) $= 7:1$.
\end{vivaqa}

\begin{vivaqa}[Question 2 (Moderate)]
Differentiate between a Plain Scale and a Diagonal Scale. State the underlying geometric principle of a Diagonal Scale.\\[4pt]
\textbf{Answer:} A Plain Scale reads two units or one unit and its first subdivision. A Diagonal Scale reads three consecutive units or up to two decimal places. Its underlying principle is the **Principle of Similar Triangles**.
\end{vivaqa}

\begin{vivaqa}[Question 3 (Exam Level)]
A distance of $5\text{ cm}$ on a map represents an actual distance of $25\text{ m}$. Calculate the $R.F.$ and determine the Length of Scale ($L.O.S.$) to measure up to $80\text{ m}$.\\[4pt]
\textbf{Answer:}  
$$R.F. = \frac{5\text{ cm}}{25\text{ m}} = \frac{5}{2500} = \mathbf{\frac{1}{500}}$$
$$L.O.S. = R.F. \times \text{Max Length} = \frac{1}{500} \times 8000\text{ cm} = \mathbf{16\text{ cm}}$$
\end{vivaqa}

\subsection{Multiple-Choice Questions (MCQs)}

1. Which pencil lead grade is recommended for drawing visible object lines?  
   A. $4\text{H}$  
   B. $6\text{H}$  
   C. $\text{H}$ or $2\text{H}$  
   D. $\text{HB}$ or $\text{B}$  
   \textbf{Answer: C} — Medium grades ($\text{H}, 2\text{H}$) are used for object lines and dimensioning.

2. What is the standard ratio of length to width for dimension arrowheads?  
   A. $1:1$  
   B. $2:1$  
   C. $3:1$  
   D. $4:1$  
   \textbf{Answer: C} — Engineering arrowheads require a $3:1$ length-to-width ratio.

3. The scale factor for a reducing scale $1:50$ means:  
   A. Drawing is 50 times larger than actual object  
   B. $1\text{ cm}$ on drawing represents $50\text{ cm}$ actual length  
   C. $50\text{ cm}$ on drawing represents $1\text{ cm}$ actual length  
   D. Representative fraction is 50  
   \textbf{Answer: B} — $R.F. = 1/50$, drawing is $1/50^{\text{th}}$ of actual size.

4. Which function key toggles Ortho mode in AutoCAD?  
   A. \texttt{F3}  
   B. \texttt{F7}  
   C. \texttt{F8}  
   D. \texttt{F12}  
   \textbf{Answer: C} — \texttt{F8} toggles Ortho mode.

5. What is the correct sequence of AutoCAD commands to set up an A4 sheet limit?  
   A. \texttt{UNITS} $\to$ \texttt{LIMITS} ($0,0$ to $210,297$) $\to$ \texttt{ZOOM} $\to$ \texttt{ALL}  
   B. \texttt{LIMITS} $\to$ \texttt{OSNAP} $\to$ \texttt{GRID}  
   C. \texttt{LINE} $\to$ \texttt{PLINE} $\to$ \texttt{RECTANG}  
   D. \texttt{DIMSTYLE} $\to$ \texttt{UNITS} $\to$ \texttt{SAVE}  
   \textbf{Answer: A} — Standard setup workflow.

---

\section{I. Unit 1 Quick Revision Checklist}
\begin{revisionchecklist}
\begin{itemize}
    \item $\text{Hard Pencils } (4\text{H--}6\text{H}) \to \text{Construction/Center Lines}$; $\text{Medium Pencils } (\text{H}, 2\text{H}) \to \text{Object Lines/Lettering}$.
    \item $\text{Visible Lines} \to \text{Continuous Thick}$; $\text{Hidden Lines} \to \text{Dashed Medium}$; $\text{Center Lines} \to \text{Chain Thin}$.
    \item $\text{Dimension Arrowhead Ratio} = 3:1$.
    \item $\text{Aligned System} \to \text{Text perpendicular to dimension line}$; $\text{Unidirectional System} \to \text{Text horizontal}$.
    \item $R.F. = \frac{\text{Drawing Length}}{\text{Actual Length}}$; $L.O.S. = R.F. \times \text{Max Length}$.
    \item $\text{Diagonal Scale Principle} = \text{Similar Triangles}$.
    \item $\text{AutoCAD Setup} \to \texttt{UNITS} \to \texttt{LIMITS } (0,0 \text{ to } 210,297) \to \texttt{ZOOM ALL}$.
    \item $\texttt{F3} = \text{OSNAP}, \quad \texttt{F8} = \text{ORTHO}, \quad \texttt{F7} = \text{GRID}$.
\end{itemize}
\end{revisionchecklist}
